% Teil II: der durchgerechnete Weg.

\subsection*{Gegeben und gesucht}
Gegeben: $T=2400$\,Jahre, $1\,\mathrm{AE}=\qty{1.5e8}{km}$,
$r_P=\qty{1}{AE}$, $v_P=\qty{42}{km/s}$, $M_\odot=\qty{1.989e30}{kg}$,
$m_\mathrm{Komet}=\qty{1.3e16}{kg}$.
Gesucht: (a) $a$; (b) $r_A$; (c) $v_P/v_A$; (d) Zeit Perihel$\to$Aphel bei
konstanter Mittengeschwindigkeit; (e) $E_\mathrm{pot}$ an beiden Punkten.

\subsection*{Ansatz}
(a) 3. Kepler in Jahren/AE. (b) $2a = r_P+r_A$. (c)
Drehimpulserhaltung. (d) Weg = halber Ellipsenumlauf $\approx \pi a$,
Zeit $=\text{Weg}/v_\mathrm{mitte}$. (e) $E_\mathrm{pot}=-GM_\odot m/r$.

\subsection*{Rechnung}
\textbf{a) Grosse Halbachse:}
\begin{equation}
  T^2 = a^3\;\Rightarrow\;
  a = T^{2/3} = 2400^{2/3}\,\mathrm{AE} \approx \qty{179.3}{AE}.
\end{equation}

\textbf{b) Aphel:}
\begin{equation}
  r_A = 2a-r_P = 2\cdot 179.3 - 1 = \qty{357.5}{AE}.
\end{equation}

\textbf{c) Geschwindigkeitsverhältnis:}
\begin{equation}
  \frac{v_P}{v_A} = \frac{r_A}{r_P}
  = \frac{357.5}{1} = 357.5
  \;(\Rightarrow v_A\approx\qty{0.117}{km/s}).
\end{equation}

\textbf{d) Zeit Perihel$\to$Aphel.} Weg$\,\approx\pi a$ (halber Umlauf).
Mittlere Geschwindigkeit $\bar v=(v_P+v_A)/2=(42+0.117)/2\approx\qty{21.06}{km/s}$:
\begin{equation}
  \pi a = \pi\cdot 179.3\cdot 1.5e8\ \mathrm{km} \approx \qty{8.44e10}{km},
\end{equation}
\begin{equation}
  \Delta t = \frac{\pi a}{\bar v}
  = \frac{8.44e10}{21.06}\ \si{s}
  \approx \qty{4.01e9}{s}
  \approx \qty{127}{a}\ \mathrm{(Jahre)}.
\end{equation}

\textbf{e) Potentielle Energien:}
\begin{equation}
  E_P = -G\frac{M_\odot m}{r_P}
  = -\frac{6.674e-11\cdot 1.989e30\cdot 1.3e16}{1.5e11}\ \si{J}
  \approx \qty{-1.15e25}{J},
\end{equation}
\begin{equation}
  E_A = -\frac{G M_\odot m}{r_A}
  = -\frac{6.674e-11\cdot 1.989e30\cdot 1.3e16}{357.5\cdot 1.5e11}\ \si{J}
  \approx \qty{-3.22e22}{J}.
\end{equation}

\begin{ergebnis}
(a) $a\approx\qty{179.3}{AE}$; (b) $r_A\approx\qty{357.5}{AE}$;
(c) $v_P/v_A=357.5$; (d) $\Delta t\approx\qty{127}{Jahre}$;
(e) $E_P\approx\qty{-1.15e25}{J}$, $E_A\approx\qty{-3.22e22}{J}$.
\end{ergebnis}

\begin{kontrolle}
\textbf{Grenzfall (c):} im Perihel ist der Komet der Sonne am nächsten und am
schnellsten; das 357-fache Geschwindigkeitsverhältnis entspricht dem
Aphel/Perihel-Verhältnis. Plausibel für einen langperiodischen Kometen.

\textbf{Einheiten:} $\Delta t$: km/(km/s)$=$s; Umrechnung in Jahre korrekt.
$E_\mathrm{pot}$: $m^3/(kg\,s^2)\cdot kg\cdot kg/m = kg\,m^2/s^2 = J$. \checkmark

\textbf{Konsistenz:} $E_P/E_A = r_A/r_P = 357.5$; Zahlen: $1.15e25/3.22e22
\approx 357$. \checkmark
\end{kontrolle}

\begin{merksatz}
Mit dem dritten Keplerschen Gesetz ($T^2=a^3$ in Jahren/AE) erhält man aus
der Periode sofort $a$; Aphel und Perihel verhalten sich bezüglich
Geschwindigkeit reziprok (Drehimpulserhaltung). Potentielle Energien
skalieren wie $1/r$.
\end{merksatz}
